% !PRINTED 580
\weberpage{580}

%<<CONT
folgt, indem wir die Ableitung nach $t$ bilden, was offenbar gestattet ist:
\[
F'(t) \equiv e\,[P(t)]^{e-1}\,P'(t)\,\Phi(t) + [P(t)]^e\,\Phi'(t)\ (\mathrm{mod}\ p).
\]
Setzen wir hierin $t=\tau$, so geht $P'(t)$ in eine durch $\mathfrak{p}$ untheilbare Form $P'(\tau)$ über (was wir schon im Beweis von \S.\,157, 4. gezeigt und benutzt haben) und wir erhalten
% [sic: Weber prints $F_1(\tau)$ here; the preceding factorization and following sentence use $\Phi(\tau)$.]
\weq{11}{F'(\tau) \equiv e\, P(\tau)^{e-1}\, P'(\tau)\, F_1(\tau)\ (\mathrm{mod}\ \mathfrak{p}^e).}
Nun ist (nach \S.\,157, 4.) $P(\tau)$ durch $\mathfrak{p}$, aber nicht durch $\mathfrak{p}^2$ theilbar, $P'(\tau)$ und $\Phi(\tau)$ sind durch $\mathfrak{p}$ nicht theilbar, und so giebt uns also die Formel (11) den Beweis des folgenden Satzes:

1. \gesperrt{Ist $\mathfrak{p}$ ein beliebiges Primideal, $p$ die durch $\mathfrak{p}$ theilbare natürliche Primzahl, und $\mathfrak{p}^e$ die höchste in $p$ aufgehende Potenz von $\mathfrak{p}$, so ist die Grundform $F'(\tau)$ allemal theilbar durch $\mathfrak{p}^{e-1}$; ist ferner der Exponent $e$ nicht theilbar durch $p$, so ist $F'(\tau)$ nicht theilbar durch $\mathfrak{p}^e$; ist aber $e$ theilbar durch $p$, so ist $F'(\tau)$ theilbar durch $\mathfrak{p}^e$ und vielleicht durch noch höhere Potenzen von $\mathfrak{p}$.}

Wenn $e$ grösser als 1 ist, so ist hiernach $F'(\tau)$ durch $\mathfrak{p}$ theilbar, und folglich ist $N_a[F'(\tau)]$ und also auch die Grundzahl $\Delta$ durch $p$ theilbar.

Wenn aber alle Primideale $\mathfrak{p}$ nur in erster Potenz in $p$ aufgehen, so ist $F'(\tau)$ relativ prim zu $p$. Zerlegt man also $F'(\tau)$ in seine Primfactoren, so kommt darunter keiner vor, dessen Norm eine Potenz von $p$ ist, folglich ist auch $N_a[F'(\tau)]$ und $\Delta$ durch $p$ nicht theilbar. Daraus folgt dann der Satz:

2. \gesperrt{Eine natürliche Primzahl $p$ ist dann und nur dann im Körper $\Omega$ durch das Quadrat eines Primfactors theilbar, wenn $p$ in der Grundzahl von $\Omega$ aufgeht.}

Aus diesen Betrachtungen können wir noch andere wichtige Schlüsse ziehen.

Wenn wir das System der linearen Gleichungen (1) in Bezug auf $\omega_1, \omega_2, \ldots, \omega_n$ auflösen, so erhalten wir Ausdrücke mit dem Nenner $U$, der, wie wir gesehen haben, eine Einheit ist. Substituirt man diese Ausdrücke in
\weq{12}{\omega = x_1\omega_1 + x_2\omega_2 + \cdots + x_n\omega_n,}
%CONT>>
